我目前为浙江理工大学物理学专业硕士研究生，预计于2027年6月毕业。硕士阶段主要围绕低维功能材料、多铁耦合及相关量子物性开展第一性原理与多尺度模拟研究。在科研过程中，我逐渐形成了从具体物理问题出发，综合运用第一性原理、有效模型、数值模拟、程序开发和数据分析建立完整证据链的研究方式。博士阶段，我希望在保留现有理论计算和编程优势的基础上，进一步向铁电/反铁电材料、外延薄膜、异质界面及极化调控方向深入，并系统补充实验研究能力。

\section{一、研究特长}

\textbf{第一性原理计算与多尺度材料模拟}是我的主要研究特长。\par
目前较为熟悉VASP、Wannier90/WannierTools、TB2J和Spirit等研究工具，能够独立开展结构优化、电子结构、自旋轨道耦合、铁电极化与翻转路径、界面结构、磁交换作用、磁各向异性、Berry曲率及相关拓扑性质等计算，并根据具体科学问题建立有效模型和开展进一步的数值模拟。

在二维多铁材料TcIrGe$_2$Se$_6$的研究中，我从结构稳定性和铁电翻转出发，进一步分析电子结构、谷自由度及微观交换机制，并利用Wannier紧束缚模型追踪Tc--Se--Ir--Se--Tc长程超超交换通道，将局域结构变化、轨道杂化和磁相互作用联系起来。相关工作已被\textit{Physical Review B}接收。这一研究使我逐渐形成了从\textbf{结构—极化—电子结构—微观相互作用—有效模型—功能响应}逐层分析材料问题的研究思路。

我的另一项优势是\textbf{科研编程和方法开发能力}。我能够使用Python、Shell/Linux进行高性能计算任务管理、批量参数扫描、数据提取、统计分析及可视化，并根据研究问题自行编写分析程序。例如，在研究长程交换机制时，常规能带和投影态密度难以定量区分具体轨道路径，我因此编写程序解析Wannier跃迁矩阵，对不同原子和轨道之间的跃迁强度进行筛选和统计，使原本定性的机制判断转化为可计算、可比较的定量证据。

此外，我也在探索\textbf{机器学习与物理模型的结合}。在相关拓扑研究中，我尝试将Sobol参数采样、拓扑不变量计算、分层机器学习与局部参数空间搜索结合，利用数据驱动方法提高高维参数空间的搜索效率，再通过严格物理计算进行验证。相关项目最终按照不同物理模块整理形成较完整的科研代码体系。这些经历使我认识到，机器学习最有价值的用途并不是替代物理模型，而是帮助研究者提高材料筛选、参数搜索和复杂数据分析的效率。

在科研工作方式上，我比较重视\textbf{结果的可靠性和可复现性}。对于U值、SOC、$k$点密度、Wannier拟合窗口、有限尺寸模型等可能影响结果的因素，我会主动进行收敛性测试和交叉验证；面对异常结果，也更倾向于重新检查结构、参数和物理假设，而不是简单接受单次计算结果。这种“发现问题—提出假设—设计验证—形成解释”的过程，是我在硕士阶段逐渐形成的主要科研习惯。

\section{二、研究兴趣}

随着硕士阶段研究的深入，我的研究兴趣逐渐从单一材料物性的预测，转向\textbf{铁电极化、结构相变、界面效应以及多自由度耦合的微观机制}。

目前我尤其关注铁电/反铁电材料中的几个问题：极化翻转过程中原子结构如何演化；不同铁电相、反铁电相和亚稳态之间如何竞争；应变、界面、缺陷和晶向如何改变极化稳定性及翻转路径；低维和外延异质结构中，局域结构变化又如何进一步影响电子和功能性质。

\newpage
我对这些问题感兴趣的重要原因在于，它们既可以从第一性原理和原子尺度模型出发进行理论研究，也能够通过真实材料制备和先进实验表征得到直接验证。这与我博士阶段希望建立的\textbf{“理论计算—材料制备—结构表征—机制分析”}研究模式高度一致。

此外，我也对复杂极化畴、极化涡旋、极化泡等非均匀极化结构感兴趣。相比均匀铁电极化，这类结构反映了晶格、电偶极、弹性场、静电作用和界面约束之间更加复杂的竞争关系，也为进一步理解铁电材料中的局域结构和功能响应提供了新的切入点。

\section{三、博士阶段研究计划}

博士阶段，我希望以\textbf{铁电/反铁电材料、外延薄膜和异质界面中的极化调控机制}为主要研究方向，并逐步建立理论与实验相结合的研究能力。

首先，在理论计算方面，我希望继续发挥第一性原理和程序开发方面的基础，从晶体结构、对称性和能量竞争关系出发，研究不同晶相、极化取向、应变状态及界面构型之间的稳定性，并利用NEB、分子动力学或有效模型分析极化翻转路径及相变机制。

其次，我希望进一步拓展多尺度模拟能力。传统DFT在时间和空间尺度上存在明显限制，因此计划学习并尝试利用第一性原理数据构建机器学习原子势，用于研究更大尺度和有限温度条件下的\textbf{极化翻转、畴壁运动、缺陷作用及结构演化}。同时，也希望探索机器学习在成分设计、结构筛选和材料性能预测中的应用，使数据方法真正服务于具体材料物理问题。

更重要的是，我希望博士阶段系统学习实验技术。硕士阶段我的研究主要以理论计算为主，而未来我希望逐步掌握\textbf{外延薄膜制备、铁电性能测试、PFM、高分辨STEM以及同步辐射结构表征}等方法，直接参与真实材料的制备和物性研究。

我希望最终形成如下研究闭环：

\begin{center}
\sffamily\small\bfseries\color{StatementTeal}
理论预测与机制分析 $\rightarrow$ 材料和结构设计 $\rightarrow$ 薄膜及异质结构制备\\[-0.1em]
$\rightarrow$ 电学和原子尺度表征 $\rightarrow$ 实验结果反馈理论模型 $\rightarrow$ 再设计与再验证
\end{center}

通过这种方式，使计算不只是对实验结果进行事后解释，而能够真正参与材料设计和实验方案制定；同时，也利用实验现象不断修正和完善理论模型。

\section{四、职业规划}

长期而言，我希望博士毕业后继续从事科研工作，并优先考虑高校、科研院所或高水平科研平台的博士后和科研岗位。

我希望未来围绕\textbf{铁电/多铁材料、低维功能材料、界面工程和极化调控}形成具有连续性的研究方向，并逐渐建立兼具\textbf{理论计算、科研编程、机器学习和实验研究}能力的个人研究特色。

我并不希望将自己局限为只熟悉某一种计算软件或实验技术的研究人员。相比具体工具本身，我更希望具备针对科学问题选择方法、开发方法并组织完整研究方案的能力。未来如果能够独立开展科研工作，我希望形成从材料设计、理论预测、程序开发，到材料制备、先进表征和物理机制分析的完整研究体系。

硕士阶段的科研经历让我逐渐确认，自己真正感兴趣的是\textbf{发现问题、建立假设、设计验证并最终理解微观机制的过程}。博士阶段，我希望进一步提升独立发现和凝练科学问题的能力，在已有计算研究基础上拓展实验技能和多尺度研究方法，最终成长为能够在理论、计算与实验之间建立联系并独立推进科研课题的研究者。